\section{Questão 13}

Suponha que temos uma moeda que não é viciada ($P(H) = P(T) = 0.5$).
Vamos criar um experimento randômico, denotando a variável $X$ como o número de caras ($H$) 
em seis lançamentos de moeda. \\

\textbf{\sffamily \color{vermelhoescuro} Observação:} note que este é um processo de Bernoulli.

\begin{multicols}{2}
\begin{itemize}
    \item [a)] Determine a faixa dinâmica de $X$.
    \item [b)] Esboce a PMF $f_X(x)$.
    \item [c)] Esboce a CDF $F_X(x)$.
    \item [d)] Encontre $P(x \geq 6)$.
    \item [e)] Encontre $P(x \leq 2)$.
    \item [f)] Encontre $P(1 \leq x \leq 5)$.
    \item [g)] Calcule $\E{X}$.
    \item [h)] Calcule $\E{X^2}$ (use LOTUS).
    \item [i)] Com os resultados de (g) e (h), calcule $\sigma_X^2$.
\end{itemize}
\end{multicols}

\respostas

\begin{itemize}
    \item [a)] 
    A quantidade de caras pode ter sido qualquer valor inteiro entre $0$ e $6$,
    logo a faixa dinâmica é $\{0, 1, 2, 3, 4, 5, 6\}$.
    
    \item [b), c)] 
    Os esboços da PMF e CDF de $X\sim\mathrm{B}(6, P(H))$ são apresentados abaixo.
    
    \begin{minipage}{.48\linewidth}
        \centering 
        \begin{tikzpicture}[> = latex, scale=1.0]
            \def\xStretch{0.95}
            \def\yStretch{1}
            \draw [->] (-\xStretch*0.5, 0) -- (\xStretch*6.5, 0) node [right] {$x$};
            \draw [->] (0, 0) -- (0, 1.5) node [above] {$f_X(x)$};
            % f_X(x) = (n k) p^k (1-p)^(n-k)
            %        = 6! / (k!(6-k)!) 0.5^6
            \draw 
                (\xStretch*0, 0) node [below] {$0$} --++ (0, \yStretch*0.015625) node (p0) {}
                (\xStretch*1, 0) node [below] {$1$} --++ (0, \yStretch*0.09375) node (p1) {}
                (\xStretch*2, 0) node [below] {$2$} --++ (0, \yStretch*0.234375) node (p2) {}
                (\xStretch*3, 0) node [below] {$3$} --++ (0, \yStretch*0.3125) node (p3) {}
                (\xStretch*4, 0) node [below] {$4$} --++ (0, \yStretch*0.234375) node (p4) {}
                (\xStretch*5, 0) node [below] {$5$} --++ (0, \yStretch*0.09375) node (p5) {}
                (\xStretch*6, 0) node [below] {$6$} --++ (0, \yStretch*0.015625) node (p6) {}
            ;
            \fill [azulunb]
                (p0) circle (0.07) node [above right] {$\frac{1}{64}$}
                (p1) circle (0.07) node [above] {$\frac{6}{64}$}
                (p2) circle (0.07) node [above] {$\frac{15}{64}$}
                (p3) circle (0.07) node [above] {$\frac{20}{64}$}
                (p4) circle (0.07) node [above] {$\frac{15}{64}$}
                (p5) circle (0.07) node [above] {$\frac{6}{64}$}
                (p6) circle (0.07) node [above] {$\frac{1}{64}$}
            ;
            \draw (3, 1.5) node {\scriptsize
                $\binom{n}{x}P(X)^x (1-P(H))^{n-x} = 
                \dfrac{6!}{x!(6-x)!} \dfrac{1}{64}$
            };
        \end{tikzpicture}
    \end{minipage}
    %
    \begin{minipage}{.48\linewidth}
        \centering 
        \begin{tikzpicture}[> = latex, scale=1.0]
            \def\xStretch{0.95}
            \def\yStretch{1}
            \draw [->] (-\xStretch*0.5, 0) -- (\xStretch*6.5, 0) node [right] {$x$};
            \draw [->] (0, 0) -- (0, 1.5) node [above] {$F_X(x)$};
            \draw 
                (\xStretch*0, 0) node [below] {$0$} --++ (0, \yStretch*0.015625) node (p0) {}
                (\xStretch*1, 0) node [below] {$1$} --++ (0, \yStretch*0.109375) node (p1) {}
                (\xStretch*2, 0) node [below] {$2$} --++ (0, \yStretch*0.34375) node (p2) {}
                (\xStretch*3, 0) node [below] {$3$} --++ (0, \yStretch*0.65625) node (p3) {}
                (\xStretch*4, 0) node [below] {$4$} --++ (0, \yStretch*0.890625) node (p4) {}
                (\xStretch*5, 0) node [below] {$5$} --++ (0, \yStretch*0.984375) node (p5) {}
                (\xStretch*6, 0) node [below] {$6$} --++ (0, \yStretch*1) node (p6) {}
            ;
            \fill [azulunb]
                (p0) circle (0.07) node [above right] {$\frac{1}{64}$}
                (p1) circle (0.07) node [above] {$\frac{7}{64}$}
                (p2) circle (0.07) node [above] {$\frac{22}{64}$}
                (p3) circle (0.07) node [above] {$\frac{42}{64}$}
                (p4) circle (0.07) node [above] {$\frac{57}{64}$}
                (p5) circle (0.07) node [above] {$\frac{63}{64}$}
                (p6) circle (0.07) node [above] {$\frac{64}{64}$}
            ;
        \end{tikzpicture}
    \end{minipage}
    
    \item [d), e), f)] 
    Pelo esboço da CDF, tem-se
    %
    \begin{align}
        P(x \geq 6) &= 1 - F_X(5) = \frac{1}{64} \\
        P(x \leq 2) &= F_X(2) = \frac{22}{64} \\
        P(1 \leq x \leq 5) &= F_X(5) - F_X(0) = \frac{62}{64}
    \end{align}
    
    \item [g)] 
    Espera-se que $\E{X}=3$ uma vez que a PMF é simétrica ao eixo $X=3$.
    O cálculo da expectância pela definição discreta confirma a suspeita:
    \begin{align}
        \E{X} 
        &= \sum_{S_X} x f_X(x) \\
        &= 
        0\cdot \frac{1}{64} +
        1\cdot \frac{6}{64} +
        2\cdot \frac{15}{64} +
        3\cdot \frac{20}{64} +
        4\cdot \frac{15}{64} +
        5\cdot \frac{6}{64} +
        6\cdot \frac{1}{64}
        =
        \frac{192}{64}
        =
        3.
    \end{align}
    
    \item [h)] 
    No caso discreto, a expectância de $Y=g(X)$ se torna, por LOTUS,
    %
    \begin{equation}
        \E{Y} = \E{g(X)} = \sum_{S_X} g(x) f_X(x).
    \end{equation}
    % 
    Especificando $g(x) = x^2$, obtém-se $\E{X^2}$:
    %
    \begin{align}
        \E{X^2} 
        &= \sum_{S_X} x^2 f_X(x) \\
        &= 
        0\cdot \frac{1}{64} +
        1\cdot \frac{6}{64} +
        4\cdot \frac{15}{64} +
        9\cdot \frac{20}{64} +
        16\cdot \frac{15}{64} +
        25\cdot \frac{6}{64} +
        36\cdot \frac{1}{64}
        =
        \frac{672}{64}
        =
        10.5.
    \end{align}
    
    \item [i)] $\sigma_X^2 = \E{X^2} - \E{X}^2 = 10.5 - 9 = 1.5$.
\end{itemize}